\documentclass[a4paper,12pt]{article}

\usepackage[T1]{fontenc}
\usepackage[utf8]{inputenc}
\usepackage[italian]{babel}
\usepackage{lmodern}
\usepackage{microtype}
\usepackage{amsmath, amssymb}
\usepackage{graphicx}
\usepackage{geometry}
\usepackage{array}
\usepackage{longtable}
\usepackage[table]{xcolor}
\usepackage{booktabs}
\usepackage{enumitem}
\usepackage{titlesec}
\usepackage{fancyhdr}
\usepackage{etoolbox}
\usepackage{xurl}
\usepackage{hyperref}

\geometry{top=2.5cm,bottom=2.7cm,left=2.6cm,right=2.6cm,headheight=16pt}

\definecolor{MainBlue}{HTML}{1F4E79}
\definecolor{AccentBlue}{HTML}{D9EAF7}
\definecolor{HeaderBlue}{HTML}{EAF2F8}
\definecolor{RuleGray}{HTML}{B7C9D6}
\definecolor{SoftGray}{HTML}{F6F8FA}

\hypersetup{
    colorlinks=true,
    linkcolor=MainBlue,
    urlcolor=MainBlue,
    citecolor=MainBlue,
    pdftitle={WP3 - Analisi di sicurezza},
    pdfauthor={Gennaro Foschillo, Francesco Vittorio Allocca}
}

\setlength{\parindent}{0pt}
\setlength{\parskip}{0.62em}
\renewcommand{\arraystretch}{1.25}
\setlength{\tabcolsep}{7pt}
\arrayrulecolor{RuleGray}
\setlength{\arrayrulewidth}{0.45pt}
\setlist[itemize]{topsep=0.25em,itemsep=0.2em,leftmargin=1.8em}
\setlist[enumerate]{topsep=0.25em,itemsep=0.2em,leftmargin=2em}
\AtBeginEnvironment{longtable}{\small}

\titleformat{\section}
  {\Large\bfseries\color{MainBlue}}
  {\thesection}{0.8em}{}
  [\vspace{0.15em}\titlerule]

\titleformat{\subsection}
  {\large\bfseries\color{MainBlue}}
  {\thesubsection}{0.8em}{}

\titleformat{\subsubsection}
  {\normalsize\bfseries\color{MainBlue}}
  {\thesubsubsection}{0.7em}{}

\titlespacing*{\section}{0pt}{1.2em}{0.8em}
\titlespacing*{\subsection}{0pt}{1.1em}{0.45em}
\titlespacing*{\subsubsection}{0pt}{0.9em}{0.25em}

\pagestyle{fancy}
\fancyhf{}
\lhead{\small Algoritmi e Protocolli per la Sicurezza}
\rhead{\small WP3 - Analisi di sicurezza}
\cfoot{\small\thepage}
\renewcommand{\headrulewidth}{0.4pt}
\renewcommand{\footrulewidth}{0pt}

\begin{document}

\pagenumbering{roman}
\begin{titlepage}
    \centering
    \vspace*{1.7cm}
    {\Large\color{MainBlue}\textbf{Algoritmi e Protocolli per la Sicurezza}\\[0.4cm]}
    {\Huge\bfseries Project Work\\[0.25cm]}
    {\LARGE\bfseries WP3: Analisi di sicurezza\\[0.6cm]}
    \textcolor{MainBlue}{\rule{0.82\textwidth}{1.2pt}}\\[1.0cm]
    {\large Analisi critica del protocollo di voto elettronico proposto nel WP2}\\[2.2cm]
    \begin{tabular}{rl}
        \textbf{Autori:} & Gennaro Foschillo \\
                         & Francesco Vittorio Allocca \\
    \end{tabular}
    \vfill
    {\large A.A. 2025--2026}
\end{titlepage}

\tableofcontents
\newpage
\pagenumbering{arabic}

\section{WP3: Analisi di sicurezza}

\subsection{Obiettivo del Work Package}

Questo Work Package analizza la sicurezza della soluzione proposta nel WP2 rispetto al modello definito nel WP1.
L'obiettivo non \`e dimostrare in modo formale e assoluto la sicurezza del sistema, ma valutare in modo critico quali propriet\`a risultano supportate dal protocollo, sotto quali assunzioni, e quali rischi residui rimangono.

L'analisi considera il protocollo per un referendum binario, basato sulla separazione tra autenticazione, autorizzazione, raccolta della scheda, registrazione pubblica e scrutinio finale.
Le propriet\`a analizzate sono autenticit\`a, unicit\`a, segretezza del voto, pseudoanonimato, integrit\`a, verificabilit\`a individuale, verificabilit\`a universale ed efficienza.

Nel seguito le propriet\`a non vengono trattate come tutte garantite allo stesso livello.
Alcune sono supportate direttamente dai meccanismi del protocollo, altre sono mitigate solo sotto specifiche assunzioni, mentre alcune rimangono rischi residui o fuori dallo scope della soluzione proposta.

\subsection{Metodo di analisi}

L'analisi viene svolta confrontando, per ciascuna propriet\`a di sicurezza:

\begin{itemize}
    \item l'obiettivo dichiarato nel modello del WP1;
    \item i meccanismi introdotti nel protocollo del WP2;
    \item le assunzioni necessarie affinch\'e tali meccanismi siano efficaci;
    \item gli scenari di attacco che il protocollo mitiga;
    \item i limiti residui che non vengono eliminati dalla soluzione proposta.
\end{itemize}

Si assume che le primitive crittografiche utilizzate siano implementate correttamente e che gli algoritmi adottati siano sicuri rispetto ad avversari computazionalmente limitati.
In particolare, l'analisi si basa su firme digitali non falsificabili, funzioni hash resistenti alle collisioni, cifratura asimmetrica probabilistica per la scheda, canali sicuri autenticati e strutture Merkle-based per la verificabilit\`a di inclusione.

\subsection{Assunzioni di sicurezza considerate}

Le propriet\`a del protocollo dipendono da alcune assunzioni operative e crittografiche.
Queste assunzioni non eliminano i rischi, ma delimitano il contesto in cui la soluzione viene valutata.

\begin{longtable}{|>{\raggedright\arraybackslash}p{0.30\textwidth}|>{\raggedright\arraybackslash}p{0.60\textwidth}|}
    \hline
    \rowcolor{HeaderBlue}\textbf{Assunzione} & \textbf{Descrizione} \\
    \hline
    \endfirsthead
    \hline
    \rowcolor{HeaderBlue}\textbf{Assunzione} & \textbf{Descrizione} \\
    \hline
    \endhead
    Primitive crittografiche sicure & Le firme digitali, gli hash, la cifratura probabilistica e i MAC usati nei canali sicuri sono considerati sicuri rispetto ad avversari efficienti. \\
    \hline
    Corretta gestione delle chiavi & Le chiavi private delle entit\`a infrastrutturali non devono essere divulgate. La compromissione di una chiave privata pu\`o invalidare le garanzie associate all'entit\`a corrispondente. \\
    \hline
    PKI e CRL disponibili & I certificati e le liste di revoca pubblicate dalla Government Authority devono essere recuperabili, firmati e aggiornati. Le entit\`a devono rifiutare certificati scaduti o revocati. \\
    \hline
    Voting Client non pienamente fidato & Il client pu\`o essere compromesso o eseguito su un dispositivo non sicuro. Il protocollo mitiga questo rischio, ma non garantisce pienamente che la preferenza mostrata e quella cifrata coincidano. \\
    \hline
    Separazione delle responsabilit\`a & La soluzione assume che le entit\`a non colludano tutte tra loro. La separazione tra GA, VS, ER e SA riduce il rischio di associazione identit\`a-voto, ma non elimina il rischio in caso di collusione estesa. \\
    \hline
    Scrutiny Authority distribuita & Almeno \(k\) nodi della Scrutiny Authority devono autorizzare lo scrutinio finale. Questa scelta \`e una custodia/autorizzazione distribuita dell'uso di \(sk_{SA}\), non una threshold decryption matematica. \\
    \hline
    Sincronizzazione temporale ragionevole & Timestamp, finestre di validit\`a, CRL e chiusura delle urne richiedono una sincronizzazione temporale sufficiente tra le entit\`a. \\
    \hline
\end{longtable}

\subsection{Classificazione delle propriet\`a}

La tabella seguente riassume il livello con cui il protocollo supporta le propriet\`a dichiarate nel modello.
La classificazione distingue tra propriet\`a supportate direttamente, propriet\`a mitigate sotto assunzioni e propriet\`a non pienamente garantite.

\begin{longtable}{|>{\raggedright\arraybackslash}p{0.20\textwidth}|>{\raggedright\arraybackslash}p{0.23\textwidth}|>{\raggedright\arraybackslash}p{0.27\textwidth}|>{\raggedright\arraybackslash}p{0.20\textwidth}|}
    \hline
    \rowcolor{HeaderBlue}\textbf{Propriet\`a} & \textbf{Livello} & \textbf{Meccanismi principali} & \textbf{Limite principale} \\
    \hline
    \endfirsthead
    \hline
    \rowcolor{HeaderBlue}\textbf{Propriet\`a} & \textbf{Livello} & \textbf{Meccanismi principali} & \textbf{Limite principale} \\
    \hline
    \endhead
    Autenticit\`a & Supportata & Autenticazione presso GA, token firmato, certificati, firma con \(sk_U\) & Dipende da GA e dalla protezione di \(sk_U\). \\
    \hline
    Unicit\`a & Supportata & \(IssuedTokensDB\), \(SpentTokensDB\), stati \emph{pending}/\emph{committed} & Dipende dalla corretta emissione dei token e dallo stato del VS. \\
    \hline
    Segretezza & Mitigata & Cifratura probabilistica con \(pk_{SA}\), separazione VS/ER/SA & Non resiste a client compromesso o collusione critica. \\
    \hline
    Pseudoanonimato & Mitigato & Separazione identit\`a-token-voto, mancata pubblicazione del token & Possibile correlazione tramite log, timestamp o ordine dei record. \\
    \hline
    Integrit\`a & Supportata & Firme digitali, MAC/canale sicuro, hash, Merkle Tree & Dipende dalla validit\`a delle chiavi e dalla correttezza del registro. \\
    \hline
    Verificabilit\`a individuale & Supportata con limite & Ricevuta firmata, \(Leaf_i\), \(MerklePath_i^{final}\), \(Root_{final}\) & Verifica l'inclusione, non il fatto che il client abbia cifrato la scelta desiderata. \\
    \hline
    Verificabilit\`a universale & Supportata se il tally pubblica dati sufficienti & \(BB_{final}\), \(Root_{final}\), preferenze decifrate validate, firme SA & Aumenta il rischio residuo di correlazione con metadati esterni. \\
    \hline
    Efficienza & Vincolo progettuale soddisfatto & Referendum binario, hash, Merkle path logaritmico, chiavi di sessione & Non misura sperimentale in WP3; sar\`a valutata in WP4. \\
    \hline
    Coercion-resistance & Non garantita & Ricevuta non contiene la preferenza & Il voto remoto non presidiato resta vulnerabile a coercizione esterna. \\
    \hline
    Disponibilit\`a & Non garantita crittograficamente & Hard-fail su CRL, separazione dei ruoli & Mancano meccanismi specifici anti-DDoS o di fault tolerance. \\
    \hline
\end{longtable}

\subsection{Analisi per propriet\`a}

\subsubsection{Autenticit\`a dell'elettore e del voto}

L'autenticit\`a richiede che soltanto elettori legittimi possano partecipare alla votazione.
Nel protocollo, questa propriet\`a \`e supportata dalla separazione tra autenticazione presso la Government Authority e successivo rilascio del token di autorizzazione.

La Government Authority autentica l'elettore e verifica la sua idoneit\`a.
Solo in caso di esito positivo viene emesso un \(VotingToken_U\), contenente un permesso \(Permit_U\) firmato dalla GA.
Il token contiene anche \(pk_U\), cio\`e la chiave pubblica associata all'elettore per la specifica elezione.
Durante l'invio della scheda, il Voting Server verifica la firma della GA sul token e la firma \(\sigma_U^{Ballot}\) prodotta con \(sk_U\).

Questa scelta impedisce a un soggetto che possiede solo il token, ma non la chiave privata corrispondente, di usarlo correttamente.
La propriet\`a dipende dalla correttezza della GA, dalla non compromissione di \(sk_U\) e dalla corretta verifica dei certificati e delle firme.

\subsubsection{Unicit\`a del voto}

L'unicit\`a richiede che ciascun elettore possa contribuire al pi\`u con una scheda valida al risultato finale.
Il protocollo affronta il problema in due punti distinti.

In fase di autorizzazione, la Government Authority aggiorna \(IssuedTokensDB\), impedendo l'emissione di pi\`u token per lo stesso elettore nella medesima elezione.
In fase di raccolta, il Voting Server controlla \(SpentTokensDB\), memorizzando un digest non invertibile del numero seriale del token.
L'uso degli stati \emph{pending} e \emph{committed} evita che un token venga considerato definitivamente speso prima che il Registro Elettorale abbia confermato la pubblicazione della scheda.

La garanzia ottenuta \`e quindi: per ogni token validamente emesso, al pi\`u una scheda pu\`o arrivare allo stato \emph{committed} ed essere associata a un record del Bulletin Board finale.
La propriet\`a non elimina il rischio di comportamento malevolo della GA, che potrebbe emettere token indebiti o duplicati; tale rischio rientra nelle assunzioni di fiducia e va discusso come rischio residuo.

\subsubsection{Segretezza della preferenza}

La segretezza della preferenza viene garantita principalmente cifrando localmente il voto con la chiave pubblica elettorale della Scrutiny Authority:

\[
    EncVoto_{SA} = Enc^{prob}_{pk_{SA}}(v).
\]

La cifratura deve essere probabilistica, ad esempio tramite RSA-OAEP, perch\'e nel caso di referendum binario una cifratura deterministica produrrebbe lo stesso ciphertext per lo stesso valore di voto.
Un osservatore del Bulletin Board potrebbe quindi raggruppare ciphertext uguali e ricavare informazioni statistiche sui voti.

Il Voting Server e il Registro Elettorale non possiedono \(sk_{SA}\) e trattano \(EncVoto_{SA}\) come oggetto opaco.
La segretezza \`e quindi preservata rispetto a tali entit\`a, sotto l'assunzione che la chiave privata di scrutinio non sia usata prima della chiusura dell'elezione e che la Scrutiny Authority segua la procedura prevista.

La propriet\`a non \`e assoluta: un Voting Client compromesso pu\`o osservare o alterare la preferenza prima della cifratura; inoltre, una compromissione della chiave privata di scrutinio o una collusione critica nella fase di scrutinio pu\`o violare la segretezza.

\subsubsection{Pseudoanonimato e separazione delle informazioni}

Il protocollo non promette anonimato perfetto, ma mira a garantire pseudoanonimato.
L'obiettivo \`e impedire che una singola entit\`a possa associare direttamente identit\`a reale, autorizzazione al voto e preferenza espressa.

La Government Authority conosce l'identit\`a reale dell'elettore e rilascia il token, ma non riceve la scheda cifrata.
Il Voting Server riceve token e scheda cifrata, ma non dovrebbe conoscere l'identit\`a reale dell'elettore n\'e la preferenza in chiaro.
Il Registro Elettorale pubblica record relativi a schede cifrate, senza token e senza identit\`a.
La Scrutiny Authority decifra le schede solo dopo la chiusura, ma non riceve informazioni di autenticazione o idoneit\`a.

Rimangono tuttavia rischi residui di correlazione.
Un avversario che combini log temporali della Government Authority, del Voting Server e del Registro Elettorale potrebbe tentare di correlare l'ordine di voto con l'identit\`a degli elettori.
Il protocollo mitiga tale rischio separando le componenti e non pubblicando i token, ma non introduce meccanismi pi\`u avanzati come mixnet o anonimizzazione forte.

\subsubsection{Integrit\`a dei messaggi e del registro}

L'integrit\`a viene supportata da firme digitali, canali sicuri, hash e struttura Merkle-based.
La cifratura della scheda protegge la segretezza, ma non viene considerata sufficiente a garantire integrit\`a.
Per questo motivo il protocollo usa firme sui principali oggetti pubblici e messaggi critici.

La Government Authority firma certificati, token, parametri pubblici e CRL.
Il Voting Server firma il messaggio inoltrato al Registro Elettorale.
Il Registro Elettorale firma ricevute e radice finale.
I nodi della Scrutiny Authority firmano le autorizzazioni alla fase di scrutinio.

Il Bulletin Board append-only, organizzato tramite foglie e radice Merkle, rende rilevabili modifiche, rimozioni o sostituzioni dei record dopo la pubblicazione.
La sicurezza di questo meccanismo dipende dalla resistenza alle collisioni della funzione hash e dalla corretta pubblicazione della radice finale firmata.

\subsubsection{Verificabilit\`a individuale}

La verificabilit\`a individuale permette all'elettore di controllare che la propria scheda cifrata sia inclusa nel registro finale.
Dopo l'accettazione della scheda, il Registro Elettorale produce una ricevuta contenente \(Leaf_i\) e la firma del Registro.
Alla chiusura dell'elezione, il Registro pubblica \(Root_{final}\) e il relativo \(MerklePath_i^{final}\).

L'elettore pu\`o quindi verificare che la foglia indicata nella ricevuta appartenga al Merkle Tree finale e che la radice sia firmata dal Registro Elettorale.
Questa verifica prova l'inclusione della scheda cifrata, non la correttezza della preferenza cifrata dal client.

Il protocollo supporta quindi una forma di verificabilit\`a individuale di inclusione.
Non garantisce pienamente la propriet\`a di \emph{cast-as-intended} in presenza di Voting Client compromesso.

\subsubsection{Verificabilit\`a universale}

La verificabilit\`a universale richiede che osservatori esterni possano controllare la coerenza tra record pubblicati, schede scrutinabili e risultato finale.
Nel protocollo, un verificatore pu\`o controllare \(BB_{final}\), \(Root_{final}\), la firma del Registro Elettorale, le autorizzazioni firmate dai nodi della Scrutiny Authority e la coerenza numerica tra schede accettate, token in stato \emph{committed} e risultato finale.

Per rendere la verificabilit\`a universale effettiva rispetto al conteggio, il \(TallyBoard\) dovrebbe pubblicare anche l'elenco delle preferenze decifrate e validate, indicizzate rispetto ai record finali:

\[
    \{(i, v_i, status_i)\}_{i=1}^{m}.
\]

In questo modo, chiunque pu\`o ricalcolare:

\[
    m = count_{\text{S\`i}} + count_{\text{No}} + count_{invalid}.
\]

Questa scelta rafforza la trasparenza del conteggio, ma introduce un compromesso: la pubblicazione dei valori decifrati associati agli indici dei record pu\`o aumentare il rischio residuo di correlazione con metadati temporali o log esterni.
Il protocollo accetta questo compromesso per privilegiare la verificabilit\`a universale richiesta dalla traccia.

\subsubsection{Efficienza}

L'efficienza viene considerata come vincolo progettuale e non solo come metrica sperimentale.
Il protocollo adotta un referendum binario, riducendo la complessit\`a della scheda e dello scrutinio.
Usa cifratura asimmetrica solo per oggetti piccoli o per proteggere la preferenza verso la Scrutiny Authority, mentre le comunicazioni operative avvengono su canali sicuri basati su chiavi di sessione.

Il controllo di riuso del token richiede una ricerca su un digest di dimensione fissa in \(SpentTokensDB\).
La verificabilit\`a individuale basata su Merkle Tree richiede prove di dimensione logaritmica rispetto al numero di schede pubblicate.
La soluzione evita inoltre blockchain pubbliche o protocolli di consenso distribuito, che aumenterebbero complessit\`a, latenza e costo computazionale.

\subsection{Analisi rispetto agli avversari del modello}

\subsubsection{Elettore malevolo}

Un elettore malevolo pu\`o tentare di votare pi\`u volte, riutilizzare un token, presentare una scheda non valida o ripudiare il proprio voto.
Il protocollo mitiga il voto multiplo tramite \(IssuedTokensDB\) e \(SpentTokensDB\).
La verifica della firma della GA sul token e della firma \(\sigma_U^{Ballot}\) impedisce l'uso di token non validi o non associati alla chiave dell'elettore.

Un elettore pu\`o comunque tentare di mostrare a terzi elementi della propria sessione di voto.
La ricevuta di inclusione non contiene la preferenza, ma il protocollo non garantisce piena resistenza alla coercizione in un contesto remoto non presidiato.

\subsubsection{Voting Server malevolo}

Un Voting Server malevolo potrebbe tentare di rifiutare schede valide, accettare schede duplicate, modificare schede cifrate o non inoltrarle al Registro Elettorale.
La modifica delle schede o dei messaggi verso il Registro \`e mitigata tramite firme e verifica del Registro Elettorale.
Il controllo \emph{pending}/\emph{committed} riduce il rischio di incoerenza tra token consumati e schede pubblicate.

Il Voting Server resta per\`o una componente critica per la disponibilit\`a e pu\`o causare denial-of-service rifiutando richieste legittime.
Questo rischio non viene eliminato crittograficamente dal protocollo.

\subsubsection{Registro Elettorale malevolo}

Un Registro Elettorale malevolo potrebbe tentare di omettere, modificare o riordinare record.
La struttura append-only e la radice Merkle firmata rendono rilevabili modifiche successive alla pubblicazione finale.
Le ricevute firmate permettono agli elettori di contestare l'eventuale mancata inclusione della propria scheda.

Il Registro conserva comunque un ruolo importante: se rifiuta di pubblicare record o non rende disponibile il Bulletin Board finale, si produce un problema di disponibilit\`a e trasparenza che deve essere trattato come rischio residuo.

\subsubsection{Government Authority malevola o compromessa}

La Government Authority \`e una componente ad alta fiducia, perch\'e autentica gli elettori, verifica l'idoneit\`a, firma token, certificati, parametri pubblici e CRL.
Una GA malevola potrebbe rilasciare token indebiti, negare autorizzazioni legittime, revocare certificati in modo scorretto o conservare log utili alla correlazione.

Il protocollo riduce la capacit\`a della GA di conoscere la preferenza, perch\'e non riceve la scheda cifrata e non partecipa allo scrutinio.
Tuttavia, non pu\`o eliminare completamente il rischio derivante da una GA malevola, che resta un'assunzione di fiducia esplicita.

\subsubsection{Scrutiny Authority collusa o compromessa}

La Scrutiny Authority decifra le schede solo dopo la chiusura delle urne.
Nel protocollo base, la distribuzione della SA non viene modellata come threshold decryption matematica, ma come procedura di custodia e autorizzazione congiunta dell'uso di \(sk_{SA}\).

Se meno di \(k\) nodi cooperano, lo scrutinio non dovrebbe essere autorizzato.
Se almeno \(k\) nodi colludono o se la procedura di custodia di \(sk_{SA}\) viene aggirata, la protezione offerta dalla distribuzione non \`e sufficiente.
Questo rischio deve essere dichiarato chiaramente: la soglia \(\langle k,n\rangle\) riduce il singolo punto di fiducia operativo, ma non fornisce una garanzia crittografica equivalente a una vera threshold cryptography.

\subsubsection{Voting Client compromesso}

Il Voting Client \`e una delle principali fonti di rischio residuo.
Un client compromesso pu\`o modificare la preferenza prima della cifratura, mostrare opzioni false, conservare dati locali o comunicare la scelta a terzi.

Il protocollo mitiga questo rischio tramite distribuzione ufficiale del client, verifica della firma del pacchetto applicativo, canali sicuri e verificabilit\`a di inclusione.
Tuttavia, non pu\`o garantire che la preferenza cifrata corrisponda a quella effettivamente voluta dall'elettore.

\subsubsection{Avversario di rete}

Un avversario di rete pu\`o osservare, intercettare, modificare o riprodurre messaggi.
Il protocollo mitiga questi attacchi tramite canali sicuri autenticati, certificati, verifica CRL, nonce, timestamp e firme sui messaggi rilevanti.

I replay vengono mitigati tramite valori freschi e controllo di stato, in particolare sul token di voto.
Rimangono possibili attacchi di disponibilit\`a, come blocco delle comunicazioni o ritardi artificiali, che non vengono risolti dalla sola crittografia.

\subsection{Compromessi progettuali}

Il protocollo adotta alcuni compromessi espliciti.
Tali compromessi non rappresentano errori della soluzione, ma scelte motivate rispetto agli obiettivi della traccia e agli strumenti trattati nel corso.

\begin{longtable}{|>{\raggedright\arraybackslash}p{0.28\textwidth}|>{\raggedright\arraybackslash}p{0.34\textwidth}|>{\raggedright\arraybackslash}p{0.28\textwidth}|}
    \hline
    \rowcolor{HeaderBlue}\textbf{Scelta} & \textbf{Motivazione} & \textbf{Rischio residuo} \\
    \hline
    \endfirsthead
    \hline
    \rowcolor{HeaderBlue}\textbf{Scelta} & \textbf{Motivazione} & \textbf{Rischio residuo} \\
    \hline
    \endhead
    Referendum binario & Riduce complessit\`a della scheda e dello scrutinio. & Minore generalit\`a rispetto a elezioni complesse. \\
    \hline
    Nessuna sostituzione del voto & Mantiene semplice l'unicit\`a del token e riduce ambiguit\`a nello scrutinio. & L'elettore non pu\`o correggere un voto inviato. \\
    \hline
    Pubblicazione delle preferenze decifrate validate & Rafforza la verificabilit\`a universale del conteggio. & Possibile correlazione con ordine o timestamp dei record. \\
    \hline
    Custodia distribuita non-threshold di \(sk_{SA}\) & Evita primitive avanzate non trattate come soluzione principale. & Non offre garanzia matematica di threshold decryption. \\
    \hline
    Nessuna blockchain pubblica & Riduce complessit\`a, latenza e overengineering. & La trasparenza dipende dal Bulletin Board e dalle firme pubblicate. \\
    \hline
    Nessuna mixnet & Mantiene il protocollo aderente agli strumenti del corso. & Resta rischio residuo di correlazione temporale. \\
    \hline
\end{longtable}

\subsection{Rischi residui}

I principali rischi residui sono i seguenti:

\begin{itemize}
    \item compromissione del Voting Client o del dispositivo dell'elettore;
    \item collusione tra pi\`u entit\`a infrastrutturali, in particolare tra GA, VS, ER e SA;
    \item uso improprio o compromissione della chiave privata di scrutinio \(sk_{SA}\);
    \item correlazione tramite log temporali, indirizzi di rete o ordine dei record pubblicati;
    \item attacchi alla disponibilit\`a di GA, VS, ER o canali di comunicazione;
    \item revoche non disponibili o CRL non aggiornate;
    \item assenza di una piena resistenza alla coercizione nel voto remoto non presidiato.
\end{itemize}

Questi rischi non invalidano il protocollo, ma delimitano il livello di sicurezza effettivamente raggiunto.
La soluzione proposta mira a un compromesso ragionevole tra sicurezza, verificabilit\`a, efficienza e semplicit\`a progettuale, senza introdurre primitive avanzate non necessarie rispetto agli obiettivi del project work.

\subsection{Possibili estensioni future non adottate}

Alcune estensioni potrebbero migliorare specifiche propriet\`a del sistema, ma non vengono adottate come parte principale del protocollo per evitare overengineering e mantenere coerenza con gli strumenti trattati nel corso.

\begin{itemize}
    \item \textbf{Threshold decryption crittografica:} migliorerebbe la protezione della chiave di scrutinio, ma richiederebbe una costruzione specifica come threshold RSA o threshold ElGamal.
    \item \textbf{Mixnet:} ridurrebbe il rischio di correlazione tra ordine di arrivo e ordine di scrutinio, ma aumenterebbe la complessit\`a del protocollo.
    \item \textbf{Cifratura omomorfa:} permetterebbe di calcolare il risultato senza decifrare singole schede, ma introduce primitive pi\`u avanzate.
    \item \textbf{Blockchain pubblica:} potrebbe aumentare la trasparenza del registro, ma introdurrebbe costi, latenza e problemi di governance non necessari per il modello scelto.
    \item \textbf{Voto sostituibile:} potrebbe mitigare coercizione e errori dell'elettore, ma complicherebbe la gestione di token, versioni della scheda e verificabilit\`a.
\end{itemize}

\subsection{Conclusione dell'analisi}

L'analisi mostra che il protocollo supporta in modo pi\`u solido autenticit\`a, unicit\`a, integrit\`a e verificabilit\`a di inclusione.
Segretezza e pseudoanonimato sono mitigati tramite cifratura, separazione delle responsabilit\`a e mancata pubblicazione dei token, ma restano dipendenti dalle assunzioni su client, Scrutiny Authority e assenza di collusione estesa.

La verificabilit\`a universale pu\`o essere resa pi\`u forte pubblicando nel \(TallyBoard\) le preferenze decifrate e validate, collegate ai record finali.
Questa scelta rende il conteggio controllabile pubblicamente, ma deve essere accompagnata dalla dichiarazione del rischio residuo di correlazione.

Il protocollo non garantisce pienamente coercion-resistance, disponibilit\`a o anonimato perfetto.
Tali limiti sono coerenti con l'obiettivo del progetto: costruire una soluzione ragionevole, efficiente e difendibile, dichiarando con precisione assunzioni e rischi residui.

\end{document}
